% 来源：OEBPS/Text/Section0001_split_003.xhtml#sigil_toc_id_28

\BookChapter{18}
\section[第18章\q 加速财富积累第2招：赚钱更多，投资更多]{第18章\\ 加速财富积累第2招：赚钱更多，投资更多}

\begin{BookQuote}
不仅要努力成为一个成功的人，而且要努力成为一个有价值的人。

\hfill ——阿尔伯特·爱因斯坦
\end{BookQuote}

好了，现在我们讲第2招。上一章讲的是省钱，还有一种比省钱更快的方式，快得几乎没有上限——只要你能释放出你的创造力和专注力，非常痴迷地找到一种方式，能为别人做得更多，比其他任何人都多，你就能赚得更多，切换到财务自由之路的快车道。

\subsection{开着卡车奔向财务自由？}

在我十几岁的时候，以后要走的路，我妈妈早就帮我想好了。她想让我当个卡车司机。她在电视上看过那些卡车驾驶技术培训学校的广告，一遍又一遍，她看上千遍也不厌倦。妈妈告诉我，只需要接受一些培训，我就能够通过考试，当上卡车司机，一年能赚2.4万美元。我的天哪，2.4万美元！我父亲在洛杉矶市区一家停车场做管理员，一年到头，辛辛苦苦才能赚1.2万美元。我开卡车一年的收入是父亲的2倍。我妈觉得，当上卡车司机，我的人生前途就会一片光明。妈妈一有机会就忽悠我：你想想看，你可以自由自在地开着你的卡车，在公路上狂奔。这确实很吸引我，至少在某种程度上如此，我想象着自己当上卡车司机，就可以打开车上的音响，一边听着我想听的音乐，一边开着车向前飞奔——对于当时才14岁大、从来没有开过车的少年来说，这真是好酷啊。当上卡车司机，我就会有机会出去闯荡，而不是像我父亲那样一辈子待在那个地下车库里，30多年，也没熬出头。

但是，后来我经历过很多悲惨的事件，我先后有过四个不同的父亲，这让我感到别人都看不起我；我从来没有足够的钱吃饱穿暖。经历的苦难多了，这也让我早熟，我意识到，我就是当上了卡车司机，开车时间再久，跑得再远，也无法逃离那种处境给我带来的痛苦。在我的头脑里，我下定决心自己以后绝对不会让我的家庭受苦受罪。除此之外，我想用我的头脑和我的心灵来赚钱，而不是当个司机靠卖力气吃饭。人生如戏，但不同的人在人生舞台上有不同的层次，我希望能够进入和父母完全不同的层次。

我看着身边其他同学和其他人的父母，心里想：别人是人，我也是人，为什么我和他们过的日子差别这么大？为什么我们家一直不断地挣扎，才能勉强吃饱肚子，才能有钱交给收账的人付清水费、电费？我们家平时吃饭只有两个选择，不是罐头装的青豆就是茄汁意大利面，为什么我们买不起番茄酱？可是就在同一个城市，住得离我们家不远，和我一起上高中的孩子，假期去的地方都是名胜古迹，上的大学的校园看起来美得如诗如画，他们过的生活明显和我过的生活不同。难道是因为他们知道我不知道的东西吗？他们做的事和我父母做的事有什么不同吗？

我开始沉溺其中。有人挣的钱是我们的2倍，但我们的工作时间是一样的，这怎么可能？有的甚至是3倍，10倍？这简直太不可思议了。从我当时的理解水平出发，这个谜题根本没法得到解答。

\subsection{投资你自己}

我的工作是个看门的保安，工资很低，因此我需要找其他活干，额外挣些钱才够吃饭。我的父母认识一个人，我父亲过去叫他“失败者”，现在这个人变得很成功，他只用了很短的时间就发达了，至少从钱上来看他是发财了。他当时正在南加利福尼亚州做房地产生意，买房，修整，装修，再转手把它高价卖出去。他需要找个男孩在周末帮他搬家具。机会来了，这个周末的短工让我乐坏了，这可能为我的人生打开一扇窗，彻底改变我的命运，我一定要好好表现，坚决不能错过。这个人的名字叫吉姆·汉纳，他肯定注意到我急匆匆地跑来跑去，非常卖力地干活。中间休息的时候，我问吉姆先生：“我想请问您，您的人生怎么会实现这么大的转变呢？您是怎么做的，才会变得这么成功呢？”

他说：“改变我人生的就是去参加了一个培训讲座，做培训讲座的人叫吉姆·罗恩。”我听不懂，问道：“什么是讲座啊？”他说：“就是找个地方，有个人来讲课，不是讲空虚无用的大道理，而是讲他过去一二十年的人生经验，他所得到的所有教训。他把所有这些经验教训的精华，浓缩成几天的课程，让你能够用几天就吸收他几十年的经验教训。”哇，这听起来简直太棒了！我也想去听。“您说的这个培训讲座，得花多少钱啊？”汉纳说：“35美元。”什么？！我干一个星期才能挣到40美元，白天读高中课程，晚上看门当保安，这40美元挣得多不容易啊。“你能带我进去听听吗？”我问他。“当然可以！”他说，“但是我不想那样做。因为你自己不花钱去听，就不会珍惜它。”我一听，愣了，站在那里，心里很难过，我怎么舍得花35美元——相当于我一个星期的工资——只是为了听这个专家讲上3个小时呢？“哦，没关系，如果你觉得花这么多钱来提升自己这项投资不值，就不要勉强去了。”他最后耸耸肩膀说道。我经过反复思想斗争，最终还是决定自己花钱去听这个培训讲座。结果这成了我一生最重要的一项投资，我花了一个星期的工资参加了这个培训讲座，遇到了讲课的老师吉姆·罗恩——这个人成了我人生的第一位导师。

我坐在加利福尼亚州欧文市一家饭店的宴会厅里，听吉姆·罗恩讲课，我完全被他迷住了。这个一头银发的男人回答的每个问题，都一直在我的大脑中盘旋。吉姆·罗恩，也和我一样，出身贫穷，曾经困惑不解，尽管他的父亲是个好人，拼命地工作，还是只能受苦受罪，而他周围的人却很成功，生活很幸福。后来，突然之间，他回答了那个问题，那个我已经问了自己好多年的问题，那个我一直找不到答案的问题。

“在经济上取得成功的秘诀是什么呢？关键是，明白自己如何在市场上变得更有价值。”

“要拥有更多，你就必须成就更多。”

“不要期望事情很容易，要期望你自己变得更好。”

“要改变事物，你要先改变自己。”

“要把事情做得更好，你要先把自己变得更好。”

我们得到报酬是因为给市场带来了价值。这需要时间……但是我们不是靠时间得到报酬的，我们得到报酬靠的是价值。美国人是独一无二的。工作报酬水平呈阶梯形，你得一步步攀登，从低工资到高工资。起点是什么样的呢？约每小时2.3美元。那么去年最高的工资收入是多少呢？那个管理迪士尼的总经理一年能挣5200万美元！这个公司怎么会付给一个人一年5200万美元呢？它傻吗？答案是这个工资收入非常合理。要是你能够帮一家公司一年赚到几十亿美元，它会付给你5200万美元吗？当然会了！这是小意思。和你创造的巨大价值相比，你拿到的这笔钱一点儿也不算多。

“你真的有可能变得那么有价值吗？答案是当然能。”接着吉姆·罗恩告诉我发财致富的终极秘密，“你如何能够真正变得更有价值，学会更加努力地提高自己的价值，比你干工作还要努力呢？”

“那么你个人能不能变得比以前的价值高2倍，用同样的时间赚的钱也能高2倍？有没有可能你个人的价值比以前高10倍，用同样的时间赚的钱也比以前高10倍？这可能吗？当然啦！”然后他停顿了一下，直视着我的眼睛，说道，“你要在同样多的时间里赚到更多的钱，你需要做的只有一件事，就是变得更有价值。”

就是这个道理！这就是我要找的答案。我知道了这个道理，它彻底改变了我的人生。就这么明白，这么简单，这句话把道理讲到我心眼儿里了——我感觉这像是当头棒喝，我突然开悟了。这句话我听吉姆·罗恩说过100多遍，一个字都没有变过。从我第一次听到之后，每一天我都把导师吉姆·罗恩的这句话牢牢记在心里，包括2009年我在他的葬礼上致辞的那一天。

就是这个人，就是这场讲座，就在这一天，吉姆·罗恩的一番话，让我重新拿回自己未来的控制权。是他让我停止专注于那些超出我控制范围的事情——我的过去，我的贫穷，其他人的预期，经济情况——是他教我专注于我能控制的东西。我可以提升自己；我可以找到办法服务别人，做得更多，变得更好，给市场增加更多的价值。我变得痴迷于一件事，找到办法为别人做得更多，比其他人做的都要多，而且用的时间更少！这是一个持续改进、持续提升的过程，永远没有终止！从本质上讲，这成了我持续进步的途径，在我做的每一个决策上，我采取的每一个行动上，它都持续不断地驱动我、引导我。

在《圣经》里，有一条简单的原则：想成为了不起的人无可厚非。如果你想成为了不起的人，你得学会成为大家的仆人，要想领导人民，你得先为人民服务。如果能找到办法，能服务很多人，你就能够赚到更多的钱。你能找到方法服务几百万人，你就能赚到几百万美元。这就是增加价值的定律：要给自己增加价值，先给别人增加价值。

可能你会觉得，巴菲特讲的比《圣经》上讲的更加容易，更通俗易懂，这个来自美国奥马哈的投资圣人有句话很有名，他说他一生做得最好的投资，也是每个人都可以做的投资，就是投资自己。巴菲特说过，他通过读提升自己能力的书来教育自己，学习戴尔·卡内基的课程完全改变了他的人生。巴菲特曾经给我讲过这个故事，那是我们俩一起在《今日秀》做节目的时候。我听了哈哈大笑，我请求他要一直不断地给大家讲这个故事。我笑着说：“讲这个对做生意很有好处。”

我把吉姆·罗恩传授给我的这个发财致富的原则牢记在心，而且一门心思地付诸行动——我从来没有停止成长，从来没有停止给予，从来没有停止努力扩大我的影响力，扩大我的能力，以便给予更多，做更多的好事。结果过了很多年之后，我在市场上变得更有价值。在这一点上，我特别幸运，现在财富再也不是我生活中的难题了。并不是我这个人独一无二，只有我才能做到，每个人都能够做到同样的事——只要你抛弃你解读自己过去经历所编的故事，突破你解读现状所编的故事，突破这些故事对你的限制。问题到处都有，办法也到处都有，只要你用心去找。

那么美国人现在的收入阶梯是什么样的呢？我敢打赌，吉姆·罗恩可能做梦也没有想到，2013年，美国收入阶梯最底层的工资收入水平是每小时7.25美元（年收入15080美元），美国当年的最高收入者是阿帕卢萨投资管理公司的创始人和对冲基金领导人戴维·泰珀，2013年个人收入达到35亿美元。有几个人能够一年赚到10亿美元呢，更别说35亿美元了！为什么有些人的收入很低，而有些人的收入高得不可思议？答案是，市场给一个麦当劳的收银员定的价值很低（每小时只有7.77美元），因为这种工作要求的技能很简单，一个人只要学几个小时就能学会。但是明显地提高大家的投资收益，这套本事就非常罕见了，当然价值也要高多了。大多数美国人都把钱存入银行，收益只有33个基点，可是戴维·泰珀能让他的投资人在同样的时间里赚到42\%的投资收益率！他对这些客户的经济贡献有多大呢？一个人能够给客户创造1\%的收益率，那他就比只会存款的普通人的价值高3倍多。戴维·泰珀能创造42\%的投资收益率，他给客户增加的经济价值就比普通人高出12627\%。

那么你要做些什么来给市场增加更多的价值呢？你怎么来保证自己能够非常富足，而不是苦苦挣扎地求生存呢？如果我们想做一个非常彻底的转变，把自己从现在所在的收入层次提高到你想要达到的更高的收入层次，实现财务自由，那么这个转变就是我所知道的能够带你到达理想之地的最有效的一条路。

可能你一听就要高喊“我反对”。别急，先听我说完。我知道，现在情况有很大不同。我知道，现在美国经济正处于困难时期。我知道，2008年以来美国损失了200多万个工作岗位，而且那些新增的工作岗位大多数都是服务业的低薪工作或者其他行业中的低薪工作。是的，我也知道从1990年到现在，美国人的平均年收入水平一直停滞不前。

你猜猜1978年，我刚开始工作时的利率水平和失业率是什么样的？在两年之内，利率飙升！我的第一笔投资是一栋四层小房子，位置在加利福尼亚州的长滩，抵押贷款利率是18\%。你能想象现在用18\%的利率来贷款买房吗？我们肯定要到美国白宫总统办公楼前抗议了。但是历史就是这样循环发展的——过去会这样，将来还会这样。但是如果你不能找到一种方法，能够有力地增加更多的价值，你的收入肯定会停滞不前。但如果你找到办法增加价值，那么你的收入变动方向只有一个，就是不断地上升，再上升。

在美国大萧条时期，5年有880万人失业。2008年金融危机，只一年就有230万人失业！那一年的失业率最高达到10\%。但是，请记住，10\%的失业率是个平均数，并不代表每个地方、每个行业都这样。部分人口的失业率超过25\%，但是对于那些年收入达到10万美元甚至更高的人来说，你猜他们的失业率是多少？答案是不到1\%。你能从中总结出什么经验教训？如果你真能训练发展自己，有一套本事是现在的市场所迫切需要的，而且你还能持续进步，变得更有价值，就会有人雇用你，或者你可以自己创业，不管经济情况怎么样，你都有工作。如果是你自己雇用自己，你越有价值，你的收入也就越高，因为你赚的钱都是你自己的。

即使今天，硅谷的就业情况也和美国其他地方的情况不同，硅谷的工作岗位比人多。那里的高科技企业发展很快，工作岗位需求大，总是招不到足够的人，因为能胜任工作的合格人才相对较少。工作岗位就摆在那里，但是你需要学习，重组自己的技能包，才能具备这些工作所需要的那些技能，才能让我们在新的市场上变得更有价值。我可以向你保证，那些“老工作”大多数一去不复返了。

我们来回顾下历史。19世纪60年代，80\%的美国人都是农民。而现在美国人里只有2\%的人在做农业工作，我们生产的粮食可以养活全世界。新的科技摧毁了一切——突然之间，一个农场工人可以干原来500个人的活儿了。很多人得很努力才能保住工作，很多人失去了原来的工作。对于那些不能适应的人来说，工业革命就是一个非常痛苦的时代。但是新科技带来了蒸汽机和其他新工具，短短的时间内就让很多人下岗了，但是它让每一个人的生活品质迅速提高，而且提供了更多收入更高的工作。

现在，新科技发展又引起了巨大的破坏，引发了大量的失业，这和上一次工业革命类似。牛津大学的研究者称，未来20年内，几乎一半的美国职业都有被自动化取代的风险（意思是机器人替代人）。我们都要重组自己的技能包，尽管每个人的重组程度会不同。我向你保证，150年前，没有一个人想过将来有一天会有这样的工作岗位：社交媒体营销人员、干细胞科学家、机器人工程师。没有一个人能想象到一个电工、一个水管工，一年能挣到15万美元；没有一个人能想象到工厂的工人能学会用电脑来实现机器自动化操作，他们一年能挣到10万美元。但是人们想象不到并不意味着这种事情不会发生。

我每天都会遇到一些人，他们告诉我就业市场比冰还要冷，公司把他们解雇了，他们担心再也找不到工作了。但是我要在这里告诉你，关键不是市场，而是你自己。你可以提高你的收入潜力，每个人都可以。你可以给市场增加价值，可以学会新的技能，可以改变你自己的思维方式。你可以成长、改变、发展，可以找到工作，获得你需要而且值得拥有的赚钱机会。

但是如果你的工作再过5\textasciitilde{}10年就变得过时了，现在你就该考虑转型，尝试一些新的东西。转型是硅谷的说法，你从一个企业换到另外一个企业，往往是在经历了一场大的惨败之后，你只能反思自己，改变自己。

你正在读本书，说明你是一个想要转变的人，正在寻找答案，寻找解决方案，寻找更好的办法。有几百种方法，你可以用来重组自己的技能包。你要做到这一点，可以去接受大学教育、商业教育、自我教育，你可以一年挣到10万美元到几百万美元，却不用辞掉工作，或者不用花一大笔钱和4年的时间只为拿个大学本科文凭（这会让你背负10万美元甚至更多的债务）。数以百万计的工作在美国各地等着人来做，但是这些工作不是随便什么人都能干的，一般人的能力都不符合工作岗位的要求。根据麦克·罗（探索频道《干尽苦差事》节目的主持人）的说法，现在美国有350万个工作岗位可供选择，其中只有10\%的岗位要求应聘者具有大学本科文凭。这意味着其他90\%的工作岗位的条件和大学文凭无关。它们更看重应聘者的培训、技能，愿意干脏活、累活，最主要的是愿意去学新的和有用的技能。正如麦克·罗所说的：“总是有工作在招人去做，但是这些工作不是我们美国传统意义上的那些工作了。”

重组自己的技能包，学习新的技能，让人既非常兴奋又非常害怕。兴奋是因为有机会去学习，去成长，去创造，去改变。兴奋的时候，你意识到：“我是有价值的，我做出了贡献，我拥有更高的身价。”你之所以害怕，是因为你会想：“我怎么能做好这个工作？”记住吉姆·罗恩的话：“对那些需要改变的事情，你必须去改变。对那些需要变得更好的事情，你必须使它们变得更好。”你要么重组自己的技能包，要么当个傻瓜让人笑话。摆脱你为解读现状所编的故事，别让这些自己编的故事限制自己，换个思路，更快地提升自己。

人们经常对我说：“托尼，要是有自己的企业，或者在一个持续成长的公司里有份工作，那当然很棒了。但是如果你做的是一份很传统的工作，工资很低，但是你又很热爱你做的这份工作，该怎么办？”我们要走出自我限制的思考方式，我来给你举一个典型的案例。那是一个老师，一度也是苦苦挣扎，他工资很低，日子过得很艰难。但是因为他对教学的热爱和他想帮助更多学生的渴望，他找到了一个方法，让他既能够为学生增加更多的价值，又能够赚到更多的钱，比大多数老师梦想的还要多。真正限制我们收入水平的，从来不是我们的工作，而是我们的创意，我们的专注，我们的贡献。

\subsection{创意、贡献和比摇滚明星还红的韩国中学老师}

你有没有遇到过一个教三年级学生的老师，他启发你尝试一些新的东西，或者遇到过一个教八年级学生的老师，他相信你的前途不可限量？你知道一个行为榜样对一个孩子一生的影响有多大吗？我们的老师就是我们最大的财富，但是老师这个职业很容易被低估，老师的薪酬也是很少的。那么，如果你是一位老师或者有份类似老师这样的低薪酬工作，上升空间有限，你能怎么办？作为一个老师，你如何给教室里30多个学生增加更多的价值？有没有一种方式，你能够给成千上万甚至上百万名学生增加更多的价值？

有很多学校的老师都会这样想：“我做我热爱的教师工作，就绝对不会赚到足够多的钱，以实现财务自由。”大家普遍认同我们这个社会赋予老师的价值低于老师按照贡献应该从社会得到的价值。但是我们现在知道，这种自我限制的想法也压抑了大多数的老师。金基勋是韩国的一位中学老师，他拒绝接受这种社会普遍认为的当老师肯定赚不了大钱的说法。

和大多数老师不同，金基勋在韩国非常有名，他就像一个摇滚明星。他是韩国最成功的老师之一。他怎么会变得如此成功？他加倍努力地提升自己，提高自己的教学能力。

60多年前，《华尔街日报》报道，韩国人大多数都是文盲。韩国政府认识到需要采取大刀阔斧的大规模行动，提高国民文化教育水平。现在，老师受到政府鼓励，他们持续不断地研究、创新。他们在教授同一个班级时，每天都会运用不同的教学方法。教育部门指导学校的老师要互相学习、互相指导——寻找最好的教学方法，来给学生增加更多的价值。结果是什么？现在韩国15岁的青少年阅读水平全球排名第二，毕业率高达93\%，而美国只有77\%。

金基勋也以优秀的老师为榜样，让自己的教学水平突飞猛进。他花了非常多的时间，分析了不同课程中最好的老师，研究了他们的教学模式，学习如何教学生能让学生的成绩实现突破。他发现了一种方法能帮助他的学生学得更快，更好，变得更聪明——他不仅教在教室里听他讲课的学生，还教整个国家的学生。他想：为什么我只专注于教教室里的30多个学生？为什么不尽我所能地帮助尽可能多的学生？利用网络科技，他意识到，他可以把自己讲的课程放到网上，让每个人都能看到他讲课的视频，让每个人都能体验到他对教学和学习的激情。

现在金基勋一周要工作60多个小时，但是其中只有3个小时花在讲课上。其余57个小时，他都用来研究、创新、开发课程，并回复学生们的提问，对他们进行辅导。“我工作越努力，我做出来的东西就越多。”他说。金基勋极其努力地工作，让自己变得更好，能为听他讲课的学生服务得更好。金基勋录下来他的讲课过程，并把它做成视频，放到网上传播，学生们可以在线听课，一小时只用花4美元。他怎么知道这样做行得通？他怎么知道自己能给学生增加比其他任何老师都更多的价值？市场总会告诉你，你的真正价值是多少，金钱不会说谎。猜一猜，有多少人买了他的课程？2013年金基勋的年收入超过400万美元，金基勋通过网络课程和网络辅导提供的价值越多，登录网站观看课程的学生人数就越多。上网听课的学生越多，就意味着金基勋赚的钱越多——不是很多，而是越来越多。

一个老师，一年能挣400万美元。和你知道的最好的老师相比，金基勋的收入要高多少呢？金基勋老师的故事突破了传统的观念：我们的职业限制了我们的收入。金基勋在韩国无疑属于金字塔顶端的那1\%的高收入人群，这并不是因为他很幸运，也不是因为他正好在合适的时间出现在了合适的地点，更不是因为他选择的职业的薪酬非常优厚。金基勋当老师能成为超级富豪，跻身韩国财富金字塔顶端1\%的人群，是因为他从来没有停止学习，从来没有停止成长，从来没有停止在自己身上投资，他不断地让自己变得更好。

\subsection{终极多任务工作者}

但是，你如果不是企业家，该怎么办？如果你根本没有兴趣干一番事业呢？如果你是在一家美国的企业工作，甚至是在一家很小的企业工作，你是不是还能找出一条路，来给别人增加更多的价值，提高你的收入潜力？我再来给你讲一个真实的故事。那是一个年轻的普通女子，名叫丹妮拉，她在一家公司的市场营销部门工作，工作岗位是艺术设计，她根本看不到一条清晰的上升通道。但是她很有才华，而且更重要的是，她渴望成长。丹妮拉总是想方设法地做更多的事情，并给予更多，这是她的天性。所以她经常主动帮助同事做一些视觉艺术设计。后来，她希望能够学到更多的市场营销知识，所以她开始研究市场营销，并主动地为同事提供帮助。后来，她发现自己对社交媒体一无所知，但是社交媒体蕴藏的市场机会潜力巨大，于是她决定研究社交媒体。她还是用同样的办法，全靠自己学，自己摸索。

几年过去了，丹妮拉一个人身兼数职，做了很多其他同事分内的工作。时间长了，这些同事忘记了丹妮拉是在帮他们干活，她是出于好心，像是送给那些同事的一份礼物。可是那些同事开始把这当作理所当然的事情，好像这活儿就应该让丹妮拉干一样。这样就形成了一个新现象，一到下午5点钟，尽管有些工作已经到了最后期限，却仍没有干完，只有丹妮拉一个人在加班，因为她的同事都从门口悄悄地溜走了。她也不想在办公室加班到很晚，但是她这个人心软，不愿意让公司和客户失望。很明显，她的那些同事其实是在占丹妮拉的便宜——丹妮拉有动力、有志向、有抱负，愿意多帮忙、多干活儿。但是丹妮拉毕竟没有三头六臂，她已经达到了极限。“我除了干好自己的活儿，还要再干三个人的活儿。”但是，她并没有抱怨，而是拿定主意，“这是个机会，我可以好好利用。”

丹妮拉怎么做的？她去找公司的首席执行官，实话实说：“现在我一个人干着四个人的活儿。我在大学里进修过很多艺术课程，在咱们公司工作这四年里，我又一边工作一边自学了视觉艺术、市场营销、社交媒体。我来找你不是要落井下石，但是我现在能给你节省一半的市场营销成本，减少三个人的工作岗位，他们三个人的活儿我一个人就能干，而且我会干得更好。光说没有用，你不需要马上信我说的话，我来证明给你看。让那三个人继续干原来的本职工作，再干六个月，我做好自己的本职工作，同时也干她们三个人的本职工作，这样你有两份成果，可以做比较，进行选择。你自己来判断是我一个干得好，还是他们三个人干得好。”

丹妮拉的要求是，6个月后，如果她一个人干的活儿比三个人加在一起干的活还要好，她的老板可以给她付双倍的工资。后来的结果是什么？丹妮拉成功了。她用工作成果证明，自己一人身兼四职完全可以胜任——视觉艺术设计出色，市场营销工作一流，文案写得很棒，社交媒体活动搞得很好。她一个人给公司增加的价值相当于4个人。公司认识到，即使给她一个人付双倍工资，公司也能够节省一半的成本。最终市场证明了一切。

\begin{BookQuote}
能够让你感到幸福的不仅是占有金钱，让你感到幸福的还有获得成就的快乐，以及创造性工作带来的兴奋。

\hfill ——富兰克林·罗斯福
\end{BookQuote}

\subsection{机会，到处都有}

你如何给这个世界增加更多的价值？你如何贡献得更多，赚得更多，影响力更大？有好几百个故事，甚至好几千个故事，讲的都是一个原来很普通的人在看到一个问题时只和其他人有些许不同，并找到了一个办法，然后采取行动，最后改变了整个行业或者创造出全新的市场。这些人并不是天生就是企业家，他们只是像你我这样的普通人，但是他们不愿意将就，不安于现状。在我们生存的这个世界中，所有的事物都是互相影响的，互联网、社交媒体、高科技，地球上的每一个人和每一件事都是互相联系的。这意味着即使是最大的、最成熟或者最稳定的企业，也会突然破产。

请在互联网搜索引擎里输入尼克·伍德曼。

\subsection{乘风破浪}

谁能预料到，20世纪主导整个摄影世界的巨无霸企业柯达公司，会在数字成像技术出现以后，走向衰败？要知道是柯达发明了数字成像技术。但是拥有124年经营历史的柯达公司在2012年宣布破产了——这个举动引发了灾难般的涟漪效应，严重影响了纽约州罗切斯特市及其周边的经济发展，超过5万人因此失业。

但是，数字成像技术这个巨大的科技和文化变革，却给一个加利福尼亚州的冲浪高手提供了巨大的机会。这个家伙叫尼克·伍德曼，狂热的冲浪爱好者。他对这项运动的纯粹热爱和投入，加上他强烈的欲望和渴求，让他能够找到一个方法来增加自我价值。

很可能你从来没有听说过伍德曼，但是他有一个非常棒的主意，他把一个防水的照相机绑在自己的腰上，这样他在冲浪的时候，能够给自己拍照。伍德曼一心想做到的就是找到一种方式，能及时地拍下他的冲浪动作。数码摄像出现之后，他又开始研究摄像机，想看看能不能把它变得更加防水，以拍摄更高质量的视频。利用新技术，他继续改进，最终他发明了极限运动的数字摄像机（GoPro），它体积很小，摄像质量却很高，可以和电视节目媲美，而且只用夹子一夹，就可以把它随身携带到任何地方。

这个很酷的小玩意，现在全世界极限运动的爱好者几乎人手一个。不管你玩的是自行车、漂流、滑板滑雪，还是冲浪，GoPro都让你能够抓住血脉偾张的神奇时刻，并分享给你所爱的每一个人。伍德曼推出GoPro运动摄像机的时机简直棒极了。他开始推销GoPro的时候，正赶上人们蜂拥地将自己的视频传到视频网站YouTube或者脸书网的时尚风潮。他创造了一个自己都很想用的产品，他觉得自己肯定不是唯一一个需要用这个东西的人。伍德曼弄清楚了如何给数百万名运动爱好者增加价值的方法，那就是让高科技更加方便、更加有趣、更加便宜。最终，伍德曼站到了时代发展潮流的前沿。这种潮流就是不管何时何地何事，人们都能积极主动地进行数字化分享。如果你真的想变得非常富有，一个核心秘密是，要站在时代发展潮流的前沿。现在这个来自加利福尼亚州圣迭戈的冲浪高手，身家超过了10亿美元。

\subsection{一个新的“类型”诞生了}

2010年，马特·劳尔邀请我参加了一场特别的圆桌会议，讨论美国经济的发展趋势。会上，我遇到了沃伦·巴菲特和全世界白手起家的最年轻的女性亿万富豪——萨拉·布雷克里。能有机会与沃伦·巴菲特一起讨论经济，已经是非常难得的了，我更没有想到的是，萨拉的故事完全让我折服了。

萨拉没有摧毁一个行业，而是创造出了一个全新的行业。她原来在迪士尼乐园工作。有一天，她准备参加派对，在换衣服时发现，自己找不到合适的内裤来搭配她那条紧身的白色裤子。她不愿意选择什么内衣也不穿，于是决定自己动手来解决问题。她手头没有缝纫工具，只有一把剪刀。于是她把连裤袜的包腿部分剪掉，瞧，一个新行业就这样诞生了。

当然，这不是一夜之间发生的事情，也不是很容易就能发生的。萨拉和我分享了一个她最重要的成功秘密：从很小的时候，她的父亲就积极地鼓励她“失败”，但是父亲定义的失败不是没能取得一个成果，而是不敢去尝试。吃晚饭的时候，父亲会问她：“今天失败过吗？”如果她今天真的失败了，也就是勇敢尝试过，父亲会相当地兴奋和激动，因为他知道这意味着孩子正在通向成功的道路上——失败是成功之母。“托尼，这治好了我害怕失败不敢尝试的毛病。”萨拉这样告诉我。

萨拉原来做过一份根本没有前途的办公用品销售工作，因此也没有很多积蓄。这一次，她拿出了自己的所有积蓄，一共5000美元，想要发明一种内衣。“我至少听到过1000次‘不行’。”她说。但她并没有听从别人的劝告。除了她投入的5000美元以外，她还省下来3000美元的律师费（其实她根本没有这么多钱）——她自己从教科书上抄了一份专利申请。

最终，她创建的公司Spanx（美国知名内衣品牌）创造了一个全新类型的产品，它被称为“塑身内衣”，引发了全世界女性的狂热追捧。按照我妻子的说法，穿上Spanx塑身内衣能把你的每一小块赘肉都装进去，你的腰围马上会小3英寸\footnote{\label{fn:pdf-18-1}1英寸＝2.54厘米。——编者注}。

著名主持人奥普拉·温弗瑞穿了这个品牌的内衣之后，在节目上说Spanx是她最爱的内衣企业。Spanx内衣马上销量大增，公司从一家小企业变成了一个轰动全球的品牌企业。现在Spanx公司的价值超过10亿美元，这个品牌旗下的产品超过200种。除了帮助女士看起来更美外，它也让她们感觉更美。萨拉天生是个乐观派，并且想把她的品牌魔力施加到我的身上：她想让我穿上一套她新推出的Spanx男士塑身内衣，当时我们俩正在《今日秀》上做节目。我向萨拉表示感谢，委婉地说，她非常了解女性内衣市场，但她可能对男性内衣市场的了解相对少一些。但萨拉勇敢创业的经历还是让我很受启发。最终Spanx男士内衣也大获成功——多亏她没有听从我的建议。现在萨拉·布雷克里百分之百地拥有Spanx公司，而且公司负债为零，从未接受过外部投资。2012年，美国《时代》杂志评选萨拉为“全世界100个最具影响力的人物”之一。

萨拉像尼克·伍德曼一样，在看到一种需求后，就采取行动满足这种需求。她拒绝接受自己过去的故事的限制，不断寻找办法来增加自我价值。

你也可以做到！你并不需要创办一家价值超过10亿美元的大公司，影响整个行业，你也不需要当个老师在网上播放讲课视频，每年赚上400万美元，你甚至不用一个人身兼数职。但是如果这些人能够想到办法赚得更多，那么你也能找到办法，每个月额外赚上500美元或者1000美元。也许，你可以一年多赚20000美元、50000美元，甚至100000美元。你能弄清楚如何释放自己内在的创造力，专注于为市场增加更多的价值，然后把你多赚的钱转入你的财务自由基金。你能做到。现在就是开始去干的时候了……

找到一种方法，多赚500美元或者多省出500美元，你一年就多6000美元。如果每年可以得到8\%的投资收益率，这样持续长期投资40年，你最后就会得到150万美元——还记得我们前面讲的不外出吃饭，订比萨外卖省钱的案例吗？如果你找到一个办法，每个月存上1000美元，相当于一年存1.2万美元，每年投资收益率为8\%，这样持续坚持40年，你最后就能积累300万美元。如果你找到一个办法，每个月存3000美元，一年就能存3.6万美元，每年投资收益率为8\%，这样持续40年，你最后就能积累900万美元的养老金。你看明白了吗？去增加价值，去赚得更多，把你赚来的钱做投资，你就可以创造出你真正渴望的那种层次的财务自由。
